%% ---------------------------------------------------------------------------
%% Buoy validation summary table (10 m wind speed + near-surface air T)
%% Pooled over the shared-year core set, PINNED to the WRF D03 column,
%% 1986--2005. Every model column pools over the identical buoy set, so the
%% pooled observation baseline is the same across all five models.
%%
%% Wind: 5 buoys (46041, 46131, 46146, 46204, 46206).
%%   (core is these 5 + 46005, but 46005 is off the D03 nest domain and is
%%   therefore pinned out of the pooled set for every column).
%% Tas: 8 buoys (46041, 46050, 46131, 46132, 46146, 46204, 46206, 46207).
%%   (core is these 8 + 46005, likewise pinned out).
%%
%% Obs: wind = 10 m anemometer magnitude (5 m vector, log-profile scaled
%%   to 10 m, z0 = 2 mm); tas = met-package air T (deg C converted to K).
%% 90 %% missing-data cascade; apples-to-apples shared-year pairing;
%% metrics on shared-year annuals (n >= 2).
%%
%% Values from results/buoy_{wind,tas}_bias_pooled.csv (relaxed gate).
%% ---------------------------------------------------------------------------
\begin{table}[t]
\centering
\caption{Pooled buoy skill against the 1986--2005 shared-year core set.
``Obs'' and ``Model'' are the shared-year period means; ``Bias'' is
Model $-$ Obs in raw units; ``Bias\,\%'' is the same bias as a fraction
of Obs; MAE and RMSE are on shared-year annuals. WRF D01/D02/D03 are the
outer (75 km)/middle (15 km)/inner (3 km) nests of the CanESM2
downscaling; CanESM2 is the parent GCM ($\sim$100 km, coarsest); CanRCM4
is a comparable $\sim$25 km downscaling. The pooled set is pinned to
the WRF D03 column, so buoy 46005 (off the D03 nest domain) is excluded
from the pooled row of every model column, giving one shared observation
baseline across all five models.}
\begin{tabular}{llrrrr}
\hline\hline
Quantity & Model & Obs & Model & Bias & Bias\,\% & MAE / RMSE$^{\dagger}$ \\
\hline
\multirow{5}{*}{\shortstack{10 m wind speed\\(m\,s$^{-1}$)}}
& CanESM2 ($\sim$100 km) & 6.00 & 5.57 & $-0.43$ & $-7.1$  & 0.81 / 0.90 \\
& WRF D01 (75 km)         & 6.00 & 5.12 & $-0.89$ & $-14.8$ & 1.51 / 1.60 \\
& WRF D02 (15 km)         & 6.00 & 6.28 & $+0.28$ & $+4.6$  & 0.71 / 0.84  \\
& WRF D03 (3 km)          & 6.00 & 6.54 & $+0.54$ & $+9.0$  & 0.66 / 0.78  \\
& CanRCM4 ($\sim$25 km)   & 6.00 & 7.00 & $+1.00$ & $+16.6$ & 1.11 / 1.23 \\
\hline
\multirow{5}{*}{\shortstack{near-surface air T\\(K)}}
& CanESM2 ($\sim$100 km) & 283.81 & 282.37 & $-1.44$ & $-0.51$ & 1.55 / 1.72 \\
& WRF D01 (75 km)         & 283.81 & 282.11 & $-1.71$ & $-0.60$ & 1.75 / 1.89 \\
& WRF D02 (15 km)         & 283.81 & 282.90 & $-0.91$ & $-0.32$ & 1.03 / 1.19 \\
& WRF D03 (3 km)          & 283.81 & 282.92 & $-0.89$ & $-0.32$ & 1.04 / 1.18 \\
& CanRCM4 ($\sim$25 km)   & 283.81 & 283.10 & $-0.71$ & $-0.25$ & 0.94 / 1.11 \\
\hline
\end{tabular}
\begin{minipage}{0.96\textwidth}
\raggedright
\footnotesize
$^{\dagger}$\,MAE / RMSE on shared-year annuals. Pooled set is pinned to
the WRF D03 column (identical across all five models): wind 5 buoys, tas
8 buoys; buoy 46005 (off the D03 nest domain) pinned out of every column.
\end{minipage}
\end{table}
